% Generated draft: requires mathematical and source review.
\documentclass[11pt]{article}
\usepackage[margin=25mm]{geometry}
\usepackage{fontspec}
\setmainfont{Latin Modern Roman}
\usepackage{amsmath,amssymb,mathtools,unicode-math}
\setmathfont{Latin Modern Math}
\usepackage{enumitem,needspace,xurl,hyperref}
\setlength{\parindent}{0pt}
\setlength{\parskip}{5pt}
\setlength{\emergencystretch}{3em}
\newcommand{\N}{\mathbb{N}}
\newcommand{\Z}{\mathbb{Z}}
\newcommand{\Q}{\mathbb{Q}}
\newcommand{\R}{\mathbb{R}}
\newcommand{\C}{\mathbb{C}}
\newcommand{\F}{\mathbb{F}}
\newcommand{\A}{\mathbb{A}}
\newcommand{\PP}{\mathbb{P}}
\newcommand{\E}{\mathbb{E}}
\newcommand{\eps}{\varepsilon}
\newcommand{\dd}{\,\mathrm{d}}
\begin{document}
\section*{UnsolvedMath ID: 700184}
\textbf{Source problem code:} catalogue entry
\subsection{Definitions and mathematical statement}
Under the usual finite-presentation interpretation of a group with a single defining relation, let \(X\) be a finite set of generators, let \(F(X)\) be the free group on \(X\), and let \(r\in F(X)\) be represented by a finite word. Define
\[
G_r=\langle X\mid r=1\rangle=F(X)/\langle\!\langle r\rangle\!\rangle,
\]
where \(\langle\!\langle r\rangle\!\rangle\) is the smallest normal subgroup of \(F(X)\) containing \(r\). For words \(u,v\) over \(X\cup X^{-1}\), write \(\bar u,\bar v\) for their images in \(G_r\). They are conjugate if
\[
\exists g\in G_r\qquad g^{-1}\bar u g=\bar v.
\]
The conjugacy problem for a fixed presentation is to find an algorithm that, on every input pair \((u,v)\), halts and correctly decides this condition.

The supplied problem, qualified as a ``well-known problem,'' asks about the conjugacy problem for groups with a single defining relation. A precise groupwise reading is: for every finite \(X\) and every relator \(r\), does \(G_r\) have a decidable conjugacy problem? Here the algorithm may depend on the presentation. A uniform reading asks whether there is one algorithm that takes \((X,r,u,v)\) as input and always halts with the correct answer. The supplied wording does not specify which quantifier order is intended, and these formulations are therefore retained separately. It also does not explicitly require finitely many generators; the finite-presentation setting above is an explicit interpretive qualification. No restriction on the relator, such as being nontrivial or not being a proper power, is stated in the record.
\subsection{Short English statement}
Can conjugacy of two words be decided algorithmically in every group defined by a single relation? The supplied wording leaves finite generation and uniformity across presentations unspecified. Its \texttt{open} label is supplied registry metadata, not an independently verified assessment.
\subsection{Sources}
\begin{itemize}
\item Supplied local UnsolvedMath registry record: permanent ID \(700184\), title \emph{Kourovka Notebook Problem 3.34}, exact source problem code \texttt{catalogue entry}, upstream revision \texttt{d371ee917ac9647fbb06cd6fc529deec6f3dda82}.
\item Supplied source attribution: \emph{The Kourovka Notebook (Unsolved Problems in Group Theory)}, 3rd issue (1969), problem 3.34; the record also identifies the 21st-issue edition (2026), \texttt{arXiv:1401.0300v48}. The supplied collection description names E. I. Khukhro and V. D. Mazurov as editors. These bibliographic details have not been independently verified.
\item The proposer field attributes the problem to D. I. Moldavanski, followed by the incomplete stored fragment \(\backslash\mathrm{u}\{\backslash\mathrm{i}\). The record gives the proposed year as 1969 and qualifies the question as a ``well-known problem.'' The malformed name ending is left unresolved.
\item The supplied registry labels the problem \texttt{open}. This label is source metadata and does not constitute an independent literature assessment. No external source was consulted for this draft.
\end{itemize}
\end{document}
